\section{CFD Model}
\label{sec:cfd_model}

\subsection{Benchmark problem}

The two-dimensional lid-driven cavity was selected as the initial benchmark
problem because it provides a simple geometry while producing a non-trivial
recirculating incompressible flow field. The problem has also been extensively
used for the verification and comparison of numerical methods. Classical
reference solutions and detailed analyses of this benchmark are available in
Ghia et al., Botella and Peyret, and Shankar and Deshpande
\cite{ghia1982,botella1998,shankar2000}.

The computational domain is a square cavity with characteristic length

\begin{equation}
L = 0.1~\mathrm{m}.
\end{equation}

The upper wall moves in the positive \(x\)-direction with constant velocity

\begin{equation}
U_{\mathrm{lid}} = 1~\mathrm{m\,s^{-1}},
\end{equation}

while the remaining cavity walls are stationary. The kinematic viscosity is

\begin{equation}
\nu = 1.0\times10^{-3}~\mathrm{m^2\,s^{-1}}.
\end{equation}

The corresponding Reynolds number is therefore

\begin{equation}
Re =
\frac{U_{\mathrm{lid}}L}{\nu}
=
100.
\end{equation}

Although the OpenFOAM computational mesh has a finite thickness in the
\(z\)-direction, the front and back boundaries are defined using the
\texttt{empty} boundary condition. The calculation therefore represents a
two-dimensional flow problem.

\subsection{Numerical solution}

The incompressible transient Navier--Stokes equations are solved using
OpenFOAM v2606. The \texttt{icoFoam} solver is used for the present benchmark.
The numerical formulation follows the finite-volume framework underlying
OpenFOAM; general finite-volume CFD methodology is described by Ferziger and
Peri\'c and by Versteeg and Malalasekera
\cite{weller1998,ferziger2002,versteeg2007,openfoam2606}.

The velocity field is denoted by \(\mathbf{U}\), and the pressure field by
\(p\). The governing equations are

\begin{equation}
\nabla\cdot\mathbf{U}=0,
\end{equation}

and

\begin{equation}
\frac{\partial\mathbf{U}}{\partial t}
+
(\mathbf{U}\cdot\nabla)\mathbf{U}
=
-\nabla p
+
\nu\nabla^2\mathbf{U}.
\end{equation}

The CFD solution obtained on the initial mesh provides the physical
information subsequently used to construct conventional and SOM-based
refinement indicators.

The machine-learning model therefore does not replace the governing
equations or the CFD solver. Instead, it operates on quantities extracted
from a physics-based CFD solution.
